\section{Conclusion}

This paper argues that proof brittleness under semantics-preserving refactoring is a measurable maintenance problem in Dafny and that this problem deserves its own benchmark-and-evaluation substrate. The main contribution is therefore not a repair algorithm, but a maintenance-centred benchmark specification for nearby-version proof breakage, together with an empirical evaluation that measures how often verification fails under controlled refactorings and how much of that failure can be recovered by simple baselines.

The intended empirical outcome is a clear residual repair gap: even after re-verification, syntactic transfer, and LLM-only repair, a substantial class of nearby-version failures should remain unresolved. By organizing those failures by refactoring class, source stratum, and proof-oriented failure mode, the paper turns an often anecdotal maintenance complaint into a reproducible benchmark problem. That is the paper's standalone value: it makes proof maintenance under evolution measurable rather than merely discussable.

\paragraph{Bridge to later work.}
The benchmark, taxonomy, and baseline results produced here also make later symbolic repair work measurable rather than merely imaginable. In that sense, this paper is both a complete empirical study and the evaluation substrate on which later work can test whether structured repair genuinely closes the residual gap identified here.
